%======================================================================
\chapter{Continuous logarithms}\label{ch:log}
%======================================================================

The basic tool of this volume is the continuous logarithm. A plane set $K$ ``goes around'' a point
$p\notin K$ exactly when the function $z\mapsto z-p$ has no continuous logarithm on $K$; the
precise forms of this statement are proved in Chapter~\ref{ch:sep}. This chapter collects the
facts about logarithms that do not involve the topology of the plane: uniqueness up to a constant,
existence on contractible spaces (from the lifting property of the exponential map), and a few
explicit logarithms. The last section proves that the plane is unicoherent and derives Brouwer's
theorem on the frontiers of complementary domains.

Throughout, points of the plane are identified with complex numbers (Section~\ref{sec:intro-conv}).

\section{Logarithms and their uniqueness}\label{sec:log-unique}

\begin{definition}\label{def:log}
Let $X$ be a topological space, $K\subseteq X$ and $f$ a complex-valued function defined (at
least) on $K$. A \emph{logarithm of $f$ on $K$} is a function $L\colon K\to\C$, continuous on $K$,
with $e^{L(z)}=f(z)$ for every $z\in K$.
\end{definition}

If $f$ has a logarithm on $K$, then $f$ vanishes nowhere on $K$, and it is continuous on $K$
(as $e^{L}$). A logarithm on $K$ restricts to a logarithm on every subset of $K$. If $L_1$ and
$L_2$ are logarithms of $f_1$ and $f_2$ on $K$, then $L_1+L_2$ is a logarithm of $f_1f_2$, $-L_1$
one of $1/f_1$, and $\cl{L_1}$ (complex conjugate) one of $\cl{f_1}$.

\begin{lemma}\label{lem:int-valued}
Let $Z$ be a preconnected subset of a topological space and $k$ a real function continuous on $Z$
whose values on $Z$ are integers. Then $k$ is constant on $Z$.
\end{lemma}

\begin{proof}
The image $k(Z)$ is a preconnected subset of $\R$, hence an interval. If it contained two integers
$m<n$ it would contain $m+\frac12$, which is not an integer.
\end{proof}

\formal{\lean{HPC.Plane.eq_of_isPreconnected_int_valued}.}

\begin{lemma}[Uniqueness]\label{lem:log-unique}
Let $S$ be a preconnected set and $L_1,L_2$ functions continuous on $S$ with $e^{L_1}=e^{L_2}$ on
$S$. Then $L_1-L_2$ is constant on $S$; the constant lies in $2\pi i\Z$.
\end{lemma}

\begin{proof}
$e^{L_1(z)-L_2(z)}=1$, so $L_1(z)-L_2(z)=2\pi i\,k(z)$ with $k(z)\in\Z$. The function
$k=\Re\big((L_1-L_2)/2\pi i\big)$ is continuous on $S$ and integer-valued, hence constant by
Lemma~\ref{lem:int-valued}.
\end{proof}

\formal{\lean{HPC.Plane.sub_const_of_exp_eq}.}

\paragraph{The principal logarithm.} Every $w\neq0$ is $w=\abs w e^{i\Arg w}$ with a unique
$\Arg w\in(-\pi,\pi]$, the \emph{principal argument}. The \emph{principal logarithm}
$\Log w=\ln\abs w+i\Arg w$ satisfies $e^{\Log w}=w$, and it is continuous on the \emph{slit plane}
$\C\setminus(-\infty,0]$. Two facts are used repeatedly:
\begin{itemize}
\item if $\abs{w-1}<1$ then $\Re w>0$, so $w$ lies in the slit plane;
\item if $\abs{u-1}<\frac12$ and $\abs{v-1}<\frac12$ then $\Log(uv)=\Log u+\Log v$. Indeed $\Re u,\Re
v>0$ and $\abs{\Arg u},\abs{\Arg v}<\pi/2$, so $\Arg u+\Arg v\in(-\pi,\pi)$ is the principal
argument of $uv$.
\end{itemize}

\formal{Mathlib's \lean{Complex.log}, \lean{Complex.arg}; the second fact is
\lean{HPC.Hagopian.log_mul_near_one}.}

\begin{lemma}[Nearby functions]\label{lem:log-near}
Let $f$ and $\Phi$ be continuous on $K$ with $\abs{\Phi(z)-f(z)}<\abs{f(z)}$ for $z\in K$. If $\Phi$
has a logarithm on $K$, so does $f$.
\end{lemma}

\begin{proof}
$f$ vanishes nowhere on $K$, and $\Phi/f=1+(\Phi-f)/f$ with $\abs{(\Phi-f)/f}<1$, so $\Phi/f$ takes
values in the slit plane and $\Log(\Phi/f)$ is continuous on $K$. If $L$ is a logarithm of $\Phi$,
then $L-\Log(\Phi/f)$ is a logarithm of $\Phi\cdot f/\Phi=f$.
\end{proof}

\formal{\lean{HPC.Moore.exists_log_of_near}.}

\section{Existence: the lifting property of the exponential}

The exponential map $\exp\colon\C\to\C\setminus\set0$ is a covering map. We use its
\emph{homotopy lifting property} without proof; it is part of the elementary theory of covering
spaces and is available in Mathlib.

\begin{theorem}[Homotopy lifting; quoted]\label{thm:lifting}
Let $X$ be a topological space, $H\colon[0,1]\times X\to\C\setminus\set0$ continuous, and
$G_0\colon X\to\C$ continuous with $e^{G_0(x)}=H(0,x)$ for all $x$. There is a continuous
$G\colon[0,1]\times X\to\C$ with $G(0,x)=G_0(x)$ and $e^{G(t,x)}=H(t,x)$ for all $t,x$.
\end{theorem}

\formal{Mathlib's \lean{IsCoveringMap.liftHomotopy} applied to \lean{Complex.isCoveringMap_exp}.}

\begin{remark}
For the reader who wants a proof: for fixed $x$, the path $t\mapsto H(t,x)$ is lifted step by
step, using that $\exp$ maps every small disc around a point of $\C$ homeomorphically onto a small
disc in $\C\setminus\set0$ (an inverse is a branch of $\Log$). Continuity in $x$ follows from the
uniqueness of lifts of paths (Lemma~\ref{lem:log-unique} on $[0,1]$) and a Lebesgue number
argument on $[0,1]\times N$ for a small neighbourhood $N$ of $x$.
\end{remark}

\begin{theorem}[Logarithms on contractible spaces]\label{thm:contractible}
Let $X$ be a topological space with a point $x_0$ and a continuous $H_c\colon[0,1]\times X\to X$
with $H_c(0,x)=x_0$ and $H_c(1,x)=x$ for all $x$. Every continuous $g\colon X\to\C\setminus\set0$
has a logarithm on $X$.
\end{theorem}

\begin{proof}
Choose $c$ with $e^c=g(x_0)$ and apply Theorem~\ref{thm:lifting} to $H(t,x)=g(H_c(t,x))$ and the
constant $G_0=c$ (note $H(0,x)=g(x_0)$). Then $x\mapsto G(1,x)$ is continuous and
$e^{G(1,x)}=g(H_c(1,x))=g(x)$.
\end{proof}

\formal{\lean{HPC.Plane.exists_exp_lift}.}

\begin{corollary}\label{cor:log-exist}
\leavevmode
\begin{enumerate}[label=(\alph*)]
\item Every continuous nowhere vanishing $g\colon\R^2\to\C$ has a logarithm on $\R^2$.
\item If $B\subseteq\R^2$ is convex and $h$ is continuous and nowhere zero on $B$, then $h$ has a
logarithm on $B$.
\item If $X$ is a topological space and $H\colon[0,1]\times X\to\C\setminus\set0$ is continuous with
$H(0,x)=1$ for all $x$, then $x\mapsto H(1,x)$ has a logarithm on $X$.
\end{enumerate}
\end{corollary}

\begin{proof}
(a) Theorem~\ref{thm:contractible} with $H_c(t,z)=tz$. (b) For $B$ nonempty pick $b_0\in B$; the
map $H_c(t,z)=b_0+t(z-b_0)$ maps $[0,1]\times B$ into $B$ by convexity, and
Theorem~\ref{thm:contractible} applies to the subspace $B$. (c) Theorem~\ref{thm:lifting} with
$G_0=0$; the logarithm is $x\mapsto G(1,x)$.
\end{proof}

\formal{(a) \lean{HPC.Moore.exists_log_pt}; (b) \lean{HPC.Moore.exists_log_of_convex}; (c)
\lean{HPC.Moore.exists_exp_lift_of_homotopy}.}

\begin{lemma}[A logarithm of $z-p$ on a small ball]\label{lem:log-ball}
Let $r>0$ and $p,c\in\R^2$ with $4r<d(c,p)$. Then $z\mapsto z-p$ has a logarithm on $B(c,r)$,
namely $L(z)=\Log\frac{z-p}{c-p}+\Log(c-p)$.
\end{lemma}

\begin{proof}
For $z\in B(c,r)$,
$\Big|\frac{z-p}{c-p}-1\Big|=\frac{\abs{z-c}}{\abs{c-p}}\le\frac{2\nrm{z-c}}{\nrm{c-p}}<\frac{2r}{4r}=\frac12$,
using $\nrm w\le\abs w\le2\nrm w$. So $(z-p)/(c-p)$ stays in the slit plane, $L$ is continuous on
the ball, and $e^{L(z)}=\frac{z-p}{c-p}(c-p)=z-p$.
\end{proof}

\formal{\lean{HPC.Plane.exists_log_on_ball}.}

\section{The plane is unicoherent; Brouwer's theorem}

Recall (\vthree{Definition~4.2}) that a connected space is \emph{unicoherent} if, whenever it is
the union of two closed connected sets, their intersection is connected. The plane is
unicoherent; the proof below is S.~Eilenberg's, by logarithms. It works for closed sets that need
not be bounded, and we state it for preconnected sets.

\begin{theorem}[Unicoherence of the plane]\label{thm:unicoherent}
Let $A,B\subseteq\R^2$ be closed and preconnected with $A\cup B=\R^2$. Then $A\cap B$ is
preconnected.
\end{theorem}

\begin{proof}
Suppose not. Then $A\cap B=P\cup Q$ with $P,Q$ disjoint, nonempty and closed (\vone{(K2)}). By
Urysohn's lemma there is a continuous $\varphi\colon\R^2\to[0,1]$ with $\varphi=0$ on $P$ and
$\varphi=1$ on $Q$. Put $e(x)=e^{2\pi i\varphi(x)}$; then $e=1$ on $A\cap B$. Define $g_0=e$ on $A$
and $g_0=1$ off $A$.

$g_0$ is continuous: both $e$ and $1$ are continuous, so it suffices that they agree on the frontier
of $A$; a frontier point of $A$ lies in $A$ (closed) and in $\cl{\R^2\setminus A}\subseteq B$, hence
in $A\cap B$, where $e=1$. Also $g_0$ never vanishes. By Corollary~\ref{cor:log-exist}(a) it has a
logarithm $G$ on $\R^2$. Put $u=\Im G/2\pi$, a continuous real function.

On $B$ we have $g_0=1$ (on $B\setminus A$ by definition, on $A\cap B$ because $e=1$ there), so
$e^{G}=1$ and $u$ is integer-valued on $B$; by Lemma~\ref{lem:int-valued} it is constant on the
preconnected set $B$. On $A$ we have $e^{G}=e^{2\pi i\varphi}$, so $u-\varphi$ is integer-valued and
hence constant on $A$. Take $p\in P$ and $q\in Q$. Both lie in $A$ and in $B$; so $u(p)=u(q)$ and
$u(p)-0=u(q)-1$, a contradiction.
\end{proof}

\formal{\lean{HPC.Moore.plane_unicoherent}; the same argument for a compact space contractible by
an explicit homotopy is \lean{HPC.Plane.isUnicoherent_univ_of_contraction}.}

\begin{lemma}\label{lem:frontier-component}
Let $W$ be an open subset of $\R^2$ and $D$ a component of $W$. Then $D$ is open and its frontier
does not meet $W$.
\end{lemma}

\begin{proof}
$D$ is open because the plane is locally connected (every point of $W$ has a ball inside $W$,
which is connected and so lies in the component of that point). If $a\in W$ were a frontier point
of $D$, the component $D'$ of $a$ in $W$ would be an open set containing $a$, hence meeting $D$; then
$D'=D$ and $a\in D$, which is impossible since $D$ is open and $a\in\cl D\setminus D$.
\end{proof}

\formal{\lean{HPC.Moore.frontier_connectedComponentIn_subset},
\lean{HPC.Plane.isOpen_connectedComponentIn_of_isOpen}.}

\begin{theorem}[Brouwer]\label{thm:brouwer}
Let $M\subseteq\R^2$ be closed and connected and $D$ a component of $\R^2\setminus M$. Then the
frontier of $D$ is connected.
\end{theorem}

\begin{proof}
$D$ is open and connected (Lemma~\ref{lem:frontier-component}).

\emph{$E=\R^2\setminus D$ is preconnected.} $E$ contains $M$. Let $y\in E\setminus M$ and let $C$ be
the component of $y$ in $\R^2\setminus M$; $C\neq D$. The closure $\cl C$ misses $D$: a point of
$D\cap\cl C$ would give, $D$ being open, a point of $D\cap C$, and then $C=D$. And $\cl C$ meets $M$:
otherwise every point of $\cl C$ lies in $\R^2\setminus M$, so by Lemma~\ref{lem:frontier-component}
it cannot be a frontier point of $C$, i.e.\ $\cl C=C$; then $C$ would be a nonempty clopen proper
subset of $\R^2$ (proper because it misses $M\neq\emptyset$). Hence $M\cup\cl C$ is connected and lies
in $E$. Every point of $E$ is thus joined to a fixed point of $M$ by a connected subset of $E$, and
$E$ is preconnected.

Now $\R^2=\cl D\cup E$ with both sets closed and preconnected, so by
Theorem~\ref{thm:unicoherent} the set $\cl D\cap E=\cl D\setminus D$ is preconnected; it is the
frontier of the open set $D$. Finally the frontier is nonempty, since otherwise $D$ would be a
nonempty clopen set missing $M$.
\end{proof}

\formal{\lean{HPC.Moore.brouwer_boundary}.}

\section*{Exercises}

\begin{exercise}
Show that $z\mapsto z$ has no logarithm on the unit circle $\set{\abs z=1}$, directly from
Lemma~\ref{lem:log-unique} applied on $[0,2\pi]$ to $t\mapsto L(e^{it})$ and $t\mapsto it$.
\end{exercise}

\begin{exercise}
Show that the closed unit disc is unicoherent, by the argument of Theorem~\ref{thm:unicoherent}.
Deduce that the circle is not a retract of the disc. (The formal development proves the
unicoherence of the disc as \lean{HPC.Plane.isUnicoherent_disk}.)
\end{exercise}

\begin{exercise}
Show that the annulus $\set{1\le\abs z\le2}$ is not unicoherent: write it as the union of its
parts in the closed upper and lower half-planes. Where does the proof of
Theorem~\ref{thm:unicoherent} break down for the annulus?
\end{exercise}
